\documentclass[11pt, a4paper]{article}

\usepackage[a4paper, top=2.5cm, bottom=2.5cm, left=2cm, right=2cm]{geometry}
\usepackage{amsmath}
\usepackage{amssymb}
\usepackage[colorlinks=true, linkcolor=blue, urlcolor=blue, citecolor=blue]{hyperref}

\title{Theory of Everything - Volume 51 \\ \Large Closed Timelike Curves (Time Travel)}
\author{Omega Protocol}
\date{\today}

\begin{document}

\maketitle

\section*{Layman's Summary}
Is time travel into the past possible? Volume 51 proves that the universe does not strictly forbid it. The geometry of spacetime can mathematically twist back on itself. However, we also prove that if you do travel back in time, the universe mathematically forces you to act in a way that prevents paradoxes.

\section{Topological Causality (Axiom 54)}
We establish that "Chronology Protection"—the idea that time must always go forward—is not an absolute, unbreakable law of physics. Instead, it is a high-probability thermodynamic limit. Because entropy usually increases, time usually flows forward. But under extreme conditions, the flow can be reversed.

\section{Closed Timelike Curves (Theorem)}
Under conditions of extreme "negative informational pressure" (such as near the core of a spinning black hole or a traversable wormhole), the network graph of the universe can form loops. We prove that these "Closed Timelike Curves" allow an observer or a piece of information to traverse a path that ends up in its own absolute past.

\section{Novikov Consistency (Theorem)}
What happens if you go back and kill your own grandfather? The Novikov Self-Consistency Principle answers this, and we prove it rigorously. The universe mathematically prohibits self-inconsistent topological loops. The probability amplitudes of a paradox sum perfectly to zero. If you travel back, the laws of physics will force you to become the cause of the past, rather than changing it.

\end{document}
